\documentclass{jsarticle}
\usepackage{ascmac,amsmath,amssymb,amsthm,bm,physics,arydshln,mathcomp,wrapfig}
\usepackage{geometry}
\usepackage[inline]{enumitem} %enumerate環境で横に箇条書きをかけるようにする
\usepackage{ulem} %波線や破線、二重線などを引けるようにする
\geometry{a4paper, left=25mm, right=25mm, top=30mm, bottom=30mm}

\begin{document}

\section*{後期・到達度評価試験に代える（大）レポート課題　（常微分方程式論２）}

\begin{flushright}
6122111　三森誠真
\end{flushright}

\begin{itembox}[l]{[1]}
区間を$I = [a,b]$ とする. $p(t) \in C^2(I), \ p(t) \neq 0$ で $r(t) \in C(I)$ とする. このとき, \ 対称作用素
$L x=p x^{\prime \prime}+p^{\prime} x^{\prime}+r x$,\\
$A x(a)+A^{\prime} x^{\prime}(a)=B x(b)+B^{\prime} x^{\prime}(b)=0$\\
を考えて, \ そのグリーン関数を $G(t, \tau, l)$ とするとき, \ 関係式
$$
\lim _{t \rightarrow \tau+0} G_{i}(t, \tau, l)-\lim _{t \rightarrow r-0} G_{t}(t, \tau, l)=\frac{1}{p(\tau)}
$$
が成り立つことを示せ. $A, \ A', \ B, \ B'$ は定数である.
\end{itembox}

\textbf{解答}．$L x=l x$ の解，$x_{1}(t, l), x_{2}(t, l)$ を，それそれ境界条件

$$
\begin{aligned}
& A x_{1}(a, l)+A^{\prime} x_{1}^{\prime}(a, l)=0 \\
& B x_{2}(b, l)+B^{\prime} x_{2}^{\prime}(b, l)=0
\end{aligned}
$$

を満たすようにとれば, 

$$
G(t, \tau, l)=\left\{\begin{array}{l}
\frac{x_{1}(\tau, l) x_{2}(t, l)}{p(\tau) W\left(x_{1}(\tau, l), x_{2}(\tau, l)\right)}, \tau \leq t \\
\frac{x_{1}(t, l) x_{2}(\tau, l)}{p(\tau) W\left(x_{1}(\tau, l), x_{2}(\tau, l)\right)}, \tau \geq t,
\end{array}\right.
$$

$W\left(x_{1}(\tau, l), x_{2}(\tau, l)\right)=x_{1}(\tau, l) x_{2}{ }^{\prime}(\tau, l)-x_{1}{ }^{\prime}(\tau, l) x_{2}(\tau, l)$.\\
ゆえに

$$
G_{t}(t, \tau, l)= \begin{cases}\frac{x_{1}(\tau, l) x_{2}^{\prime}(t, l)}{p(\tau) W\left(x_{1}(\tau, l), x_{2}(\tau, l)\right)}, & \tau<t, \\ \frac{x_{1}^{\prime}(t, l) x_{2}(\tau, l)}{p(\tau) W\left(x_{1}(\tau, l), x_{2}(\tau, l)\right)}, & \tau>t\end{cases}
$$

ゆえに

$$
\begin{gathered}
\lim _{t \rightarrow+0} G_{t}(t, \tau, l)-\lim _{t \rightarrow-0} G_{t}(t, \tau, l) \\
=\frac{x_{1}(\tau, l) x_{2}^{\prime}(\tau, l)-x_{1}^{\prime}(\tau, l) x_{2}(\tau, l)}{p(\tau) W\left(x_{1}(\tau, l), x_{2}(\tau, l)\right)}=\frac{1}{p(\tau)} .
\end{gathered}
$$

\bigskip
\bigskip

\begin{itembox}[l]{[2]}
関数のある族$\Gamma$ を 
$$
\Gamma := \{x\equiv x(t) | x \in C^2[0,1], \ かつx(0) = x(1) = 0 \}
$$
と定める. 

(1) \ 
$\Gamma$ において定義された対称作用素$L x=\frac{d^{2} x}{d t^{2}}$ の固有値およびグリーン関数を求めよ.

(2) \ 
そのグリーン関数を用いて微分方程式
$$
\frac{d^{2} x}{d t^{2}}-l x = 1
$$
の解$x(t) \in \Gamma$ を求めよ. ただし, \ $l$ は$L$ の固有値ではないとする.
\end{itembox}

\textbf{解答}
$\frac{d^{2} x}{d t^{2}} = lx$ の一般解 は

$$
x = c_{1} \cos \sqrt{-l}t+c_{\mathrm{s}} \sin \sqrt{-l}t
$$

である．$x(0)=0$ から $c_{1}=0. \ x(1)=0$ から
$c_{2} \sin \sqrt{-1}=0$

ゆえに $\sqrt{-l}=n \pi(n= \pm 1, \pm 2, \cdots)$ ，すなわち
\underline{$l = -n^2\pi^2$ が固有値である}.\\
グリーン関数を求めるために，まず $x(0)=0$ となるような $L x=l x$ の解を求めると

$$
x=x_{1}(t, l)=\sin \sqrt{-l} t
$$

また $x(1)=0$ となる解を求めると

$$
\begin{aligned}
x=x_{2}(t, l) & =\sin \sqrt{-l} \cos \sqrt{-l}t -\cos \sqrt{-l} \sin \sqrt{-l}t \\
&= \sin \sqrt{-l}(1-t) .
\end{aligned}
$$

また

$$
\begin{aligned}
& W\left(x_{1}(\tau, l), x_{2}(\tau, l)\right) \\
&= \sin \sqrt{-l}\tau(-\sqrt{-l} \cos \sqrt{-l)(1-\tau)-\sqrt{-l} \cos \sqrt{-l}\tau} \times \sin \sqrt{-l}(1-\tau) \\
&=-\sqrt{-l} \sin \sqrt{-l}(\tau+1-\tau)\\
&=-\sqrt{-l} \sin \sqrt{-l .} .
\end{aligned}
$$

ゆえに
$$
G(t, \tau, l) 
=\begin{cases}\frac{\sin{ \sqrt{-l}} \tau \sin \sqrt{l}(1-t)}{\sqrt{-l} \sin \sqrt{-l}}, & \tau \leq t, \\ 
\frac{\sin \sqrt{l} t \sin \sqrt{-l}(1-\tau)}{\sqrt{-l} \sin \sqrt{-l}} & \tau \geq t .\end{cases}
$$

$l>0$ のときは，関係
$$
\sin \sqrt{-l} t=i \sinh \sqrt{l} t
$$

を使えば，上記は次のように書き直せる.

$$
G(t, \tau, l) 
=\begin{cases}-\frac{\sinh \sqrt{l} \tau \sinh \sqrt{l}(1-t)}{\sqrt{l} \sinh \sqrt{l}}, & \tau \leq t, \\ 
-\frac{\sinh \sqrt{l} t \sinh \sqrt{l}(1-\tau)}{\sqrt{l} \sinh \sqrt{l}}, & \tau \geq t .\end{cases}
$$

次に

$$
\frac{d^{2} x}{d t^{2}}-l x=1
$$

の解で $x(t) \in \Gamma な る も の は, l$ が固有値でなければ

$$
\begin{aligned}
x
&=\int_{0}^{1} G(t, \tau, l) d \tau \\
&= -\int_0^t \frac{\sin \sqrt{-l} \tau \sin\sqrt{-l}(1-t)}{\sqrt{-l}\sin\sqrt{-l}} d\tau
- \int_t^1 \frac{\sin \sqrt{-l}t \sin\sqrt{-l}(1-\tau)}{\sqrt{-l}\sin\sqrt{-l}} d\tau \\
&= \frac{1}{\sqrt{-l} \sin \sqrt{-l}} \left[\frac{1}{\sqrt{-l}} \cos\sqrt{-l}\tau\sin\sqrt{-l}(1-t) \right]_0^t
 + \frac{1}{\sqrt{-l} \sin \sqrt{-l}} \left[-\frac{1}{\sqrt{-l}}\sin\sqrt{-l}t\cos\sqrt{-l}(1-\tau) \right]_t^1 \\
&= \frac{1}{-l\sin\sqrt{-l}} \{\cos\sqrt{-l}t\sin\sqrt{-l}(1-t)-\sin\sqrt{-l}(1-t)-\sin\sqrt{-l}t + \sin\sqrt{-l}t\cos\sqrt{-l}(1-t) \} \\
&=\frac{1}{-l \sin \sqrt{-l}} \{\sin \sqrt{-l}(t+1-t)-\sin \sqrt{-l} t-\sin \sqrt{-l}(1-t) \} \\
&= \underline{-\frac{1}{l} + \frac{1}{l\sin\sqrt{-l}}(\sin\sqrt{-l}t + \sin\sqrt{-l}(1-t))}.
\end{aligned}
$$

\bigskip
\bigskip

\begin{itembox}[l]{[3]}
閉区間[a,b] 上で定義された関数の族$\Gamma$ を
$$
\Gamma := \{x\equiv x(t) | x \in C^2[a,b], \ かつx(a) = x(b) = 0 \}
$$
と定める. この$\Gamma$ において定義された対称2階線形微分作用素
$$
Lx = px'' + p'x' + rx
$$
を考える. $L$ の係数$p, r$ は上の[1] と同じとする.

(1) \ 
$p(t)>0(a \leq t \leq b)$ ならば, \ $L$ の固有値は上に有界であることを示せ.

(2) \ 
$p(t)<0(a \leq t \leq b)$ ならば, \ $L$ の固有値は下に有界であることを示せ.
  
\end{itembox}

\textbf{解答}

固有値の1つを$\lambda$ ，それに属する固有関数を $\varphi(t)$ とする。 ただし 
$$
\int_{a}^{b}(\varphi(t))^{2} d t=1
$$

としておく.

$
L \varphi-\lambda \varphi=0$ から
$\int_{a}^{b} \varphi L \varphi d t-\lambda \int_{a}^{b} \varphi^{2} d t=0 .$

ここで $\int_{a}^{b} \varphi^{2} d t=1$ に注意すれば

$$
\begin{aligned}
\lambda= & \int_{a}^{b} \varphi L \varphi d t \\
= & \int_{a}^{b} \varphi\left(p \varphi^{\prime \prime}+p^{\prime} \varphi^{\prime}+r \varphi\right) d t \\
= & \int_{a}^{b} \varphi\left(p \varphi^{\prime}\right)^{\prime} d t+\int_{a}^{b} r \varphi^{2} d t \\
&= \left[\varphi p \varphi^{\prime}\right]_{a}^{b}-\int_{a}^{b} p \varphi^{\prime 2} d t+\int_{a}^{b} r \varphi^{3} d t .
\end{aligned}
$$

境界条件により $\left[\varphi p \varphi^{\prime}\right]_{a}^{b}=0$ ．ゆえに

$$
\lambda=-\int_{a}^{b} p \varphi^{\prime 2} d t+\int_{a}^{b} r \varphi^{2} d t .
$$

rは $[a, b]$ で連続だから有界である. そこで $m<r<M$とすれば $\int_{0}^{b} \varphi^{2} d t=1$ であるから

$$
m<\int_{a}^{b} r \varphi^{2} d t<M
$$

ゆえに

$$
m-\int_{a}^{b} p \varphi^{\prime 2} d t<\lambda<M-\int_{a}^{b} p \varphi^{\prime 2} d t
$$

$\varphi^{\prime 2} \geq 0$ だから，$p>0$ ならば $\lambda<M であり, \ p<0$ ならば $m<\lambda$ ．\\

\bigskip
\bigskip

\begin{itembox}[l]{[4]}
$C=C[a, b]$ とし，


\begin{gather*}
\Gamma=\left\{f(t) \in C^{2}[a, b] \mid A f(a)+A^{\prime} f^{\prime}(a)=B f(b)+B^{\prime} f^{\prime}(b)=0,\right. \\
\left.A^{2}+A^{\prime 2} \neq 0, B^{2}+{B^{\prime}}^{2} \neq 0\right\} \tag{8}
\end{gather*}


とする．$L$ を $\Gamma$ 上で定義された対称微分作用素


\begin{equation*}
L x=p x^{\prime \prime}+p^{\prime} x^{\prime}+r x \tag{9}
\end{equation*}


とし，$R_{\ell}$ を $L$ のレゾルベントとする. $\left\{\mu_{k}\right\}_{k=0,1,2, \ldots}$ を $R_{\ell}$ の固有値，それに対応する規格化された固有関数を $\left\{\varphi_{k}\right\}_{k=0,1,2, \ldots}$ とする. このとき関数列 $\left\{\varphi_{k}\right\}_{k}$ は完全であること を示せ. すなわち，関係式


\begin{equation*}
\sum_{k=0}^{\infty}\left|\left(f, \varphi_{k}\right)\right|^{2}=\|f\|^{2}, \quad(\forall f \in C) \tag{10}
\end{equation*}


が成り立つことを示せ.（この関係式（10）をパーセバルの関係式ということもある．）\\
  
\end{itembox}
\textbf{解答}

(i) \ $f \in \Gamma,\left(f, \varphi_{k}\right)=c_{k}$ のときを考える. 級数$\sum_{k=0}^{n} c_{k} \varphi_{k}$ は $n \rightarrow \infty$ のとぎ$f$に一様収束することがわかっている. ゆえに任意の$\varepsilon > 0$ に対してある自然数$N$ が存在し, \ $n \geq N$ のとき
$$
|f(t) - \displaystyle\sum_{k=0}^n c_k \varphi_k (t)| < \varepsilon
$$
が成り立つ. 

よって, \ $\displaystyle\lim_{n \to \infty} \left\| f - \sum_{k=0}^n c_k \varphi_k \right \|^2 = 0$ であり, \ 

$\displaystyle \left\| f - \sum_{k=0}^n c_k \varphi_k \right\|^2 = \|f\|^2 - 2\sum_{k=0}^n c_k \varphi_k (t) + \sum_{k,l = 0}^n c_k c_l (\varphi_k, \varphi_l), \ (f, \varphi_k) = c_k$.

また$\{\varphi_k\}$ は正規直交系であるから
$$
(\varphi_k, \varphi_l) = \begin{cases} 0 & (k \neq l)\\
1 & (k=l) 
\end{cases}
$$

よって, \ $\displaystyle \left\| f - \sum_{k=0}^n c_k \varphi_k \right\|^2 = \|f\|^2 - \sum_{k = 0}^n c_k^2$

ここで$n \to \infty$ とすると左辺は$0$ に収束するから, \ 
$$
\|f\|^2 = \displaystyle\lim_{n \to \infty} \sum_{k=0}^n c_k^2 = \sum_{k=0}^{\infty} (f, \varphi_k)^2.
$$

ゆえに $f \in \Gamma$ ならば，問題の結論が成り立つ.

\bigskip

(ii)

次に, \ $f \in C$ で必ずしも$f \in \Gamma$ ではない状況を考える. このとき, \ 任意の$\varepsilon > 0$ に対して $\|f - \bar{f}\| < \varepsilon$ であるような $\bar{f} \in \Gamma$ が存在することから, \ 

$$
\begin{aligned}
\| f-\sum_{k=0}^{n}\left(f, \varphi_{k}\right) \varphi_{k} \|
&= \| f- \bar{f} + \bar{f} -\sum_{k=0}^{n}\left(f, \varphi_{k}\right) \varphi_{k} +\sum_{k=0}^{n}\left(f, \varphi_{k}\right) \varphi_{k}-\sum_{k=0}^{n}\left(f, \varphi_{k}\right) \varphi_{k} \| \\
&\leq \|f-\bar{f} \|+\left\|\bar{f}-\sum_{k=0}^{n}\left(\bar{f}, \varphi_{k}\right) \varphi_{k}\right\| + \left\|\sum_{k=0}^{n}\left(f-f, \varphi_{k}\right) \varphi_{k}\right\|
\end{aligned}
$$


仮定により $\|f-\bar{f}\|<\varepsilon ───\textcircled{\scriptsize 1}$ 

また， $\bar{f} \in \Gamma$ であるから, \ $n$ を十分大きくとりさえすれば

$$
\left|f-\sum_{k=0}^{n}\left(f, \varphi_{k}\right) \varphi_{k}\right|<\varepsilon ───\textcircled{\scriptsize 2}
$$

また $\left(\varphi_{k}\right)$ の正規直交性より

$$
\begin{aligned}
\left|\sum_{k=0}^{n}\left(f-\bar{f}, \varphi_{k}\right) \varphi_{k}\right|^{2}
&= \sum_{k=0}^{n}\left(f-\bar{f}, \varphi_{k}\right)^{2} \\
&\leq\|f-\bar{f}\|^{2} \ （ベッセルの不等式）\\
&< \varepsilon^{2}
\end{aligned}
$$


ゆえに

$$
|\sum (f - \bar{f}, \varphi_k)\varphi_k| < \varepsilon ───\textcircled{\scriptsize 3}
$$

\textcircled{\scriptsize 1}, \textcircled{\scriptsize 2}, \textcircled{\scriptsize 3} より, \ $\| f-\sum_{k=0}^{n}\left(f, \varphi_{k}\right) \varphi_{k} \| < 3\epsilon$.

よって, \ $\displaystyle\lim_{n \to \infty} \| f-\sum_{k=0}^{n}\left(f, \varphi_{k}\right) \varphi_{k} \| = 0$.

したがって, \ (i) と同様にして$$
\|f\|^{2}= \sum_{k=0}^{\infty}\left(f, \varphi_{k}\right)^{2}
$$

が得られる.\\


\bigskip
\bigskip

\begin{itembox}[l]{[5]}
$C=C[a, b]$ とする．$C$ 上で定義された線形作用素 $T$ は，（i）$\|T\|<\infty$ であるとき，有界な作用素であるとよばれ，（ii） $\lim _{n \rightarrow \infty}\left\|x_{n}-x\right\|=0,\left(x_{n} \in C, x \in C\right)$ のとき，


\begin{equation*}
\lim _{n \rightarrow \infty}\left\|T x_{n}-T x\right\|=0 \tag{11}
\end{equation*}


ならば，連続な作用素であるとよばれる. このとき，作用素 $T$ の有界性と連続性とは全く同値であることを示せ.  
\end{itembox}

\textbf{解答}

（有界性$\Longrightarrow$ 連続性）

$T$ が有界であると仮定すると, \ 
$$
\| Tx_n - Tx\| = \| T(x_n - x)\| \leq \|T\| \|x_n - x\|.
$$

$\|T\| < \infty$ だから, \ $\|x_n - x\| \to 0$ ならば $\|Tx_n - Tx\| \to 0$. よって, \ $T$ は連続である.

\bigskip

（連続性$\Longrightarrow$ 有界性）

$T$ が連続であると仮定し, \ $\|y\| < 1 (y \in \mathbb{C})$ とする. このとき, \ $T$ の連続性から, \ $\|\delta y\| < \delta$ ならば $\|T\delta y\| = \delta\|Ty\| < 1$ となるような$\delta > 0$ が存在する. よって, \ $\|f\| < 1$ ならば $\|Ty\| < \frac{1}{\delta}$ すなわち $\|T\| < \infty$ であることが示された.

\bigskip
\bigskip

\begin{itembox}[l]{[6]}
関数空間 $C, \Gamma$ および $\Gamma$ 上の対称微分作用素 $L$ を問題［4］と同様に定義し，$L$ のレ ゾルベントを $R_{\ell}$ とする．$R_{\ell}$ のベキを


\begin{equation*}
R_{\ell}^{k} u=R_{\ell}\left(R_{\ell}^{k-1} u\right), \quad k=2,3, \ldots, \quad u \in C \tag{12}
\end{equation*}


で定義し，


\begin{align*}
G_{k}(t, \tau, \ell) & =\int_{a}^{b} G(t, s, \ell) G_{k-1}(s, \tau, \ell) d s, \quad k=2,3, \ldots  \tag{13}\\
G_{1}(t, \tau, \ell) & =G(t, \tau, \ell) \tag{14}
\end{align*}


とする．このとき，


\begin{equation*}
R_{\ell}^{k} u(t)=\int_{a}^{b} G_{k}(t, \tau, \ell) u(\tau) d \tau \tag{15}
\end{equation*}


であることを示せ．ただし，$G(t, \tau, \ell)$ は作用素 $L$ のグリーン関数である．  
\end{itembox}

\textbf{解答}

数学的帰納法を用いて示す.

$G_{1}(t, \tau, l) = G(t, \tau, l)$ は定義から明らか.

$k = n$ のとき
$$
R_{i}^{k} u(t)-\int_{a}^{b} G_{k}(t, \tau, l) u(\tau) d \tau,
$$
$$
G_k(t, \tau, l) = \int_a^b G(t,s,l)G_{k-1}(s,\tau,l) ds
$$

であると仮定すると, \ 

$$
\begin{aligned}
R_{l}^{a+1} u(t) & =\int_{a}^{b} G(t, s, l) R_{l}^{n} u(s) d s \\
& =\int_{a}^{b} G(t, s, l) d s \int_{a}^{b} G_{n}(s, \tau, l) u(\tau) d \tau \\
& =\int_{a}^{b}\left(\int_{a}^{b} G(t, s, l) G_{n}(s, \tau, l) d s\right) u(\tau) d \tau
\end{aligned}
$$

となり, 

$$
R_{l}^{n+1} u(t)=\int_{a}^{b} G_{n+1}(t, \tau, l) u(\tau) d \tau,
$$
$$
G_{n+1}(t, \tau, l) = \int_{a}^{b} G(t, s, l) G_{n}(s, \tau, t) d s
$$

\bigskip
\bigskip

\begin{itembox}[l]{[7]}
上記問題［6］と同じ設定とする．このとき，展開式

\begin{equation*}
G_{k}(t, \tau, \ell)=\sum_{j=1}^{\infty} \mu_{j}^{k} \varphi_{j}(t) \varphi_{j}(\tau), \quad k=2,3, \ldots \tag{16}
\end{equation*}


が成り立つことを示せ．ただし $, \mu_{0}, \mu_{1}, \mu_{2}, \ldots$ は $R_{\ell}$ の固有値で，$\varphi_{0}, \varphi_{1}, \varphi_{2}, \ldots$ はそれ に対応する規格化された固有関数である.  
\end{itembox}

\textbf{解答}

$$
G_{k}(t, \tau, D)-\int_{a}^{b} G(t, s, l) G_{k-1}(s, \tau, 1) d s
$$

であるから, \ 

$G_{k}$ を$t$ の関数と考えれば
$G_k = R_l G_{k-1}$.

よって$G_k \in \Gamma$ であり, \ $G_k(t,\tau,l) = \displaystyle \sum_{j=0}^{\infty} c_j \varphi_j(t)$
$$
\begin{aligned}
\qquad c_{j}=\left(G_{k}, \varphi_{j}\right)
&=\int_{a}^{b} G_{k}(t, \tau, l) \varphi_{j}(t) d t\\
&=\int_{a}^{b} G_{k}(\tau, t, l) \varphi_{j}(t) d t\\
&=\left(R_{l}{ }^{k} \varphi_{j}\right)(\tau)
\end{aligned}
$$

$R_{l} \varphi_{j}=\mu_{j} \varphi_{j}, \quad R_{l}{ }^{2} \varphi_{j}=R_{l}\left(\mu_{j} \varphi_{j}\right)=\mu_{j}^{2} \varphi_{j}, \ \cdots, \ R_{l}{ }^{k} \varphi_{j}=\mu_{j}^{k} \varphi_{j}$.
なので, \ $c_j = \mu_j^k \varphi_j (\tau)$. したがって

$$
G_{k}(t, \tau, l)=\sum_{j=0}^{\infty} \mu_{j}^{k} \varphi_j(t) \varphi_{j}(\tau)
$$

\bigskip
\bigskip

\begin{itembox}[l]{[8]}
関数の族 $\Gamma$ を


\begin{equation*}
\Gamma:=\left\{x \equiv x(t) \mid x \in C^{2}[0,1], x(0)=x^{\prime}(1)=0\right\} \tag{17}
\end{equation*}


で定める．$\Gamma$ を定義域にもつ作用素 $L$ を


\begin{equation*}
L x=\frac{d^{2} x}{d t^{2}} \tag{18}
\end{equation*}


とする．境界値問題 $L x=-t$ を解け．  
\end{itembox}

\textbf{解答}


$$
L x=\frac{d^{2} x}{d t^{2}}=0
$$

を解くと $x=a t+b$（ $a, b$ は定数）となる. ここで境界条件$x(0)-x^{\prime}(1)=0$ を考えると$a = b = 0$，すなわち

$$
x \equiv 0 .
$$

ゆえに$0$ は固有値ではないから$Lx = -t$ はグリーン関数を用いて解くことができる.

$$
\frac{d^{2} x}{d t^{2}}=0
$$ の解で, \ $x(0) = 0$ を満たすものとして $x_{2}(t)=1$ をとれば $W\left(x_{1}(\tau), x_{2}(\tau)\right)=-1$ たからら

$$
G(t, \tau, 0)= \begin{cases}-x_{1}(\tau) x_{2}(t)=-\tau, & \tau \leq t, \\ -x_{1}(t) x_{2}(\tau)=-t, & \tau \geqq t,\end{cases}
$$

ゆえに求める解は

$$
\begin{aligned}
x & =\int_{0}^{1} G(t, \tau, 0) \cdot(-\tau) d \tau \\
& =\int_{0}^{t} \tau^{2} d \tau+\int_{t}^{1} t \tau d \tau=\frac{t}{2}-\frac{t^{3}}{6}
\end{aligned}
$$

\bigskip
\bigskip

\begin{itembox}[l]{[9]}
関数の族 $\Gamma$ を


\begin{equation*}
\Gamma:=\left\{x \equiv x(t) \mid x \in C^{2}[0,1], x^{\prime}(0)=x^{\prime}(1)=0\right\} \tag{19}
\end{equation*}


とする．$u(t) \in C[0,1]$ に対して，$\Gamma$ において境界値問題


\begin{equation*}
L x=\frac{d^{2} x}{d t^{2}}=u(t) \tag{20}
\end{equation*}


を解け.  
\end{itembox}

\textbf{解答}

$x=1$ は $x^{\prime}(0)=x^{\prime}(1)=0$ を満たし, \ 
$\frac{d^{2} x}{d t^{2}}=0$

の解だから， $0$ は $L$ の固有値である. ゆえにこれはグリーン関数を使って解くことができない.\\

$$
\frac{d^{2} x}{d t^{2}} = l x
$$

の一般解は

$$
x=c_{1} \cos \sqrt{-l}t + c_{2} \sin \sqrt{-l} t
$$

ゆえに

$$
x' = -\sqrt{-l}c_1\sin\sqrt{-l}t + \sqrt{-l}c_2\cos\sqrt{-l}t.
$$

$x'(0) = 0$ から $c_{2}=0$ または，$l=0, x^{\prime}(1)=0$ から $l \neq 0$ ならば $\sin \sqrt{-l} = 0$ ．\\
ゆえに $\sqrt{-l} = n \pi$, すなわち $l=-n^{2} \pi^{2}$（ $n$ は整数）.\\
したがって $L$ の固有値は

$$
l--n^{2} \pi^{2} \quad(n=0,1,2, \cdots)
$$

であり，それに対応する規格化された固有関数は\\
$l=0$ に対しては $x_{0} = 1$ ，

$$
l=-n^{2} \pi^{2} \quad(n \neq 0) \text { に対しては } x_{n}=\sqrt{2} \cos n \pi t .
$$

ゆえに

$$
\frac{d^{2} x}{d t^{2}}=u(t)
$$

が $\Gamma$ の中で解をもつのは，$u(t)$ が固有値 0 に対応する固有関数 1 と直交し ているとき，すなわち

$$
\int_{0}^{1} u(t) d t=0
$$

のときに限る.\\
そしてこの時の解は

$$
\begin{aligned}
x & =\sum_{n=1}^{\infty} \frac{\left(x_{n}, u\right)}{-n^{2} \pi^{2}-0} x_{n}(t)+c x_{0}(t) \\
& =\underline{-\sum_{n=1}^{\infty} \frac{2}{n^{2} \pi^{2}} \cos n \pi t \cdot \int_{0}^{1} u(\tau) \cos n \pi \tau d \tau+c}
\end{aligned}
$$

このフーリエ数数が一様収束することに注意して答えを書き直すと,

$$
\begin{aligned}
x & =-\int_{0}^{1} u(\tau) d \tau\left[\sum_{n=1}^{\infty} \frac{2}{n^{2} \pi^{2}} \cos n \pi t \cdot \cos {n\pi \tau} \right]+c \\
& =-\int_{0}^{1} u(\tau) d \tau \left[\sum_{n=1}^{\infty} \frac{\cos n \pi(\tau+t)}{n^{2} \pi^{2}}+\sum_{n=1}^{\infty} \frac{\cos n\pi(\tau - t)}{n^2\pi^2} \right] + c
\end{aligned}
$$

次に，$[-1,1]$ において $s^{2}$ に等しく，あとはそれを周期$2$ をくりかえすような周期関数を $f(s)$ とすれば，それは

$$
f(s)=\sum_{n=0}^{\infty} f_{n} \cos n \pi s
$$

のようなフーリエ級数に展開されて，

$$
f_{0}=\int_{0}^{1} s^{2} d s=\frac{1}{3},
$$

$$
f_{n}=2 \int_{0}^{1} s^{2} \cos n \pi s d s=\frac{4(-1)^{n}}{n^{2} \pi^{2}}, \quad n=1,2, \cdots
$$

すなわち

$$
f(s)=\frac{1}{3}+4 \sum_{n=1}^{\infty} \frac{(-1)^{n}}{n^{2} \pi^{2}} \cos n \pi s
$$

また，$[-1,1]$ で $|s|$ に等しく，あとはそれを周期 $2$ をもってくりかえすような周期関数を $g(s)$ とおけば，それは

$$
g(s)=\sum_{n=0}^{\infty} g_{n} \cos n \pi s
$$

のようなフーリエ数数に展開されて

$$
g_{0}=\int_{0}^{1} s d s=\frac{1}{2}
$$

$$
g_n = 2\int_0^1 s\cos n\pi s ds 
= \begin{cases}
0 & n:偶数\\
-\frac{4}{n^2\pi^2} & n:奇数
\end{cases}
$$

すなわち $g(s) = \frac{1}{2} - 4\sum' \frac{1}{n^2\pi^2} \cos n\pi s$

ただし$\sum'$ は $n$ が奇数についてのみの和を表すものとする.

これから$f(s)-2 g(s)=-\displaystyle \frac{2}{3}+4 \sum_{n=1}^{\infty} \frac{\cos n \pi s}{n^{2} \pi^{2}}$.

したがって

$$
\begin{aligned}
x &= -\int_{0}^{1} u(\tau)\left(\frac{1}{6}+\frac{1}{4} f(\tau+t)-\frac{1}{2} g(\tau+t)+\frac{1}{6}+\frac{1}{4} f(\tau-t)-\frac{1}{2} g(\tau-t)\right) d \tau+c \\
&=-\int_{0}^{1} u(\tau)\left[\frac{1}{3}+\frac{1}{4}(f(\tau+t)-2 g(\tau+t)) + \frac{1}{4}(f(\tau-t)-2 g(\tau-t))\right] d \tau+c . \\
& f(\tau+t)= \begin{cases}(\tau+t)^{2}, & 0 \leqq \tau \leqq 1-t, \\
(\tau+t-2)^{2}, & 1-t \leqq \tau \leqq 1,\end{cases} \\
& f(\tau-t)=(\tau-t)^{2}, 0 \leqq \tau \leqq 1, \\
& g(\tau+t)= \begin{cases}\tau+t, & 0 \leqq \tau \leqq 1-t, \\
2-\tau-t, & 1-t \leqq \tau \leqq 1,\end{cases} \\
& g(\tau-t)= \begin{cases}-\tau+t, & 0 \leqq \tau \leqq t, \\
\tau-t, & t \leqq \tau \leqq 1 .\end{cases}
\end{aligned}
$$

これから

$$
\frac{1}{4}(f(\tau+t)-2 g(\tau+t))=\frac{1}{4}\left(\tau^{2}+t^{2}+2 \tau t-2 \tau-2 t\right) \ (0 \leq \tau \leq 1)
$$

$$
\frac{1}{4}(f(\tau-t)-2 g(\tau-t))=\left\{\begin{array}{l}
\frac{1}{4}\left(\tau^{2}+t^{2}-2 \tau t+2 \tau-2 t\right), \ 0 \leq \tau \leq t\\
\frac{1}{4}\left(\tau^{2}+t^{2}-2 \tau t-2 \tau+2 t\right), \ t \leq \tau \leq 1
\end{array}\right.
$$

よって,

$$
\begin{aligned}
& \frac{1}{4}(f(\tau+t)-2 g(\tau+t))+\frac{1}{4}(f(\tau-t)-2 g(\tau-t)) \\
& = \begin{cases}\frac{1}{2}\left(\tau^{2}+t^{2}\right)-t, & 0 \leq \tau \leq t, \\
\frac{1}{2}\left(\tau^{2}+t^{2}\right)-\tau, & t \leq \tau \leq 1 .\end{cases}
\end{aligned}
$$

ゆえに，解は

$$
\begin{aligned}
x= & -\frac{1}{3} \int_{0}^{1} u(\tau) d \tau-\frac{1}{2} \int_{0}^{1} \tau^{2} u(\tau) d \tau-\frac{1}{2} t^{2} \int_{0}^{1} u(\tau) d \tau \\
& +t \int_{0}^{t} u(\tau) d \tau+\int_{t}^{1} \tau u(\tau) d \tau+c
\end{aligned}
$$

$\int_{0}^{1} u(\tau) d \tau=0, \int_{0}^{1} \tau^{2} u(\tau) d \tau$ は定数であることに注意すれば，

$$
\underline{x=t \int_{0}^{t} u(\tau) d \tau+\int_{t}^{1} \tau u(\tau) d \tau+c}
$$

が求める解である.\\

\bigskip
\bigskip

\begin{itembox}[l]{[10]}
境界値問題


\begin{equation*}
\frac{d^{2} x}{d t^{2}}-\ell x=1, \quad x^{\prime}(0)=x(1)=0 \tag{21}
\end{equation*}


を解け.  
\end{itembox}

\textbf{解答}

まず
$$
\frac{d^{2} x}{d t^{2}}-l x=0
$$

の一般解は

$$
x=c_{1} \cos \sqrt{-l}t + c_{2} \sin \sqrt{-l}t .
$$

$x^{\prime}(0)=x(1)=0$ を満たすためには $c_{2}=0, \ \cos \sqrt{-l}=0$. ゆえに固有値は

$$
l = -\left(\frac{1}{2}+n\right)^{2} \pi^{2}, \quad n=0,1,2, \cdots
$$

それに対する規格化された固有関数は

$$
\sqrt{2} \cos \left(\frac{1}{2}+n\right) \pi t
$$

である．\\
したがって, 

(i) \ $l$ が固有値$-\left(\frac{1}{2}+n\right)^{2} \pi^{2}$ に等しいときは, \ 
$\int_0^1 \cos\left(\frac{1}{2} + n \right) \pi t$ \ dt = $\frac{1}{(\frac{1}{2} + n)\pi} \neq 0$ なので境界値問題は解けない.

(ii) \ $l \neq -\left(\frac{1}{2} + n \right)^2 n^2$ のときはグリーン関数を作ることができる.

まず, \ $x'(0) = 0$ となる解は $x = \varphi_1(t) = \cos\sqrt{-l}t$.

$x(1) = 0$ となる解は
$$
x = \varphi_2(t) = \sin\sqrt{-l}(t-1),
$$
$$
W(\varphi_1(t), \varphi_2(t)) = \sqrt{-l}\cos\sqrt{-l}.
$$

ゆえに \ 
$G(t, \tau, l) = \begin{cases}
\frac{\varphi_1(\tau)\varphi_2(t)}{\sqrt{-l}\cos\sqrt{-l}} = \frac{\cos\sqrt{-l}\tau\sin\sqrt{-l}(t-1)}{\sqrt{-l}\cos\sqrt{-l}}, & \tau \le t \\
\frac{\varphi_1(t)\varphi_2(\tau)}{\sqrt{-l}\cos\sqrt{-l}} = \frac{\cos\sqrt{-l}t\sin\sqrt{-l}(\tau - 1)}{\sqrt{-l}\cos\sqrt{-l}}, & \tau \ge t
\end{cases}$

したがって求める解は

$$
\begin{aligned}
x(t)
&= \int_0^1 G(t, \tau, l) d\tau \\
&= \int_0^t \frac{\cos\sqrt{-l}\tau\sin\sqrt{-l}(t-1)}{\sqrt{-l}\cos\sqrt{-l}} d\tau + \int_t^1 \frac{\cos\sqrt{-l}t\sin\sqrt{-l}(\tau - 1)}{\sqrt{-l}\cos\sqrt{-l}} d\tau \\
&=\frac{\cos\sqrt{-l}t\sin\sqrt{-l}(t - 1)}{\sqrt{-l}\cos\sqrt{-l}} \\
&= \underline{-\frac{1}{l} + \frac{\cos\sqrt{-l}t}{l\cos\sqrt{-l}}}
\end{aligned}
$$

\bigskip
\bigskip

\begin{itembox}[l]{[11]}
$\Gamma$ を問題［8］で定義された関数の族と同じものとする．$\Gamma$ 上で定義された 2 階対称線形微分作用素 $L$ のグリーン関数を $G(t, \tau, \ell)$ とするとき，


\begin{equation*}
G(t, \tau, \ell)-G(t, \tau, 0)=\ell \int_{a}^{b} G(t, s, 0) G(s, \tau, \ell) d s \tag{22}
\end{equation*}


が成り立つことを示せ．ただし， 0 と $\ell$ は共に $L$ の固有値ではないものとする．  
\end{itembox}

\textbf{解答}

$Lx = px'' + p'x' + rx$ とすれば, \ $[a, b]$ 内の任意の区間$[\alpha, \beta]$ において$2$ 回連続微分可能な関数$x(s), y(s)$ に対して

$$
\displaystyle\int_\alpha^\beta (Lx)(s)y(s) ds - \int_\alpha^\beta (Ly)(s)x(s) ds
= [p(s)(x'(s)y(s) - x(s)y'(s))]_\alpha^\beta
$$

が成り立つ.

いま, \ $x(s) = G(s,\tau, l), \ y(s) = G(s, t, 0)$ とする. まず$\tau < t, \ a < \tau_- < \tau < \tau_+ < t_- < t < t_+ < b$ とすれば, \ $x(s), y(s)$ は$[a, \tau_-], \ [\tau_+, t_-], \ [t_+, b]$ においてそれぞれ$2$ 回連続微分可能であって, \ そこで$(L-l)x(s) = 0, \ Ly(s) = 0$ が成り立つ. ゆえに, \ 上の$3$つの区間のうちの$1$ つを$[\alpha, \beta]$ と書けば

$$
\begin{aligned}
\displaystyle \int_\alpha^\beta (Lx)(s)y(s) ds - \int_\alpha^\beta (Ly)(s)x(s) ds
&= l\int_\alpha^\beta x(s)y(s) ds \\
&= l\int_\alpha^\beta G(s, \tau, l)G(s, t, 0) ds
\end{aligned}
$$

さらに
$$
\begin{aligned}
\displaystyle \int_\alpha^\beta (Lx)(s)y(s) ds - \int_\alpha^\beta (Ly)(s)x(s) ds
&= [p(s)(x'(s)y(s) - x(s)y'(s))]_\alpha^\beta \\
&= [p(s)(G_s(s, \tau, l)G(s, t, 0) - G(s, \tau, l)G_s(s, t, 0))]_\alpha^\beta  
\end{aligned}
$$

よって, \ 
$l \left\{\displaystyle\int_a^{\tau_-} G(s, \tau, l)G(s, t, 0) ds + \int_{\tau_+}^{l_-} G(s, \tau, l)G(s, t, 0) ds + \int_{l_+}^b G(s, \tau, l)G(s, t, 0) ds \right\} \\
= [p(s)(G_s(s, \tau, l)G(s, t, 0) - G(s, \tau, l)G_s(s, t, 0))]_a^{\tau_-} + [p(s)(G_s(s, \tau, l)G(s, t, 0) - G(s, \tau, l)G_s(s, t, 0))]_{\tau_+}^{t_-} + [p(s)(G_s(s, \tau, l)G(s, t, 0) - G(s, \tau, l)G_s(s, t, 0))]_{l_+}^b$.

ここで, \ $\tau_- \to \tau, \ \tau_+ \to \tau, \ t_- \to t, \ t_+ \to t$ とすれば, \ 上式の左辺は$l\displaystyle\int_a^b G(s,\tau,l)G(s,t,0) ds$ に収束する.

これに対して右辺は

$p(b)(G_s (b, \tau, l)G(b, t, 0) - G(b, \tau, l)G_s(b,t,0)) \\
- p(a)(G_s(a, \tau, l)G(a, t, 0) - G(a, \tau, l)G_s(a,t,0)) \\
- \{p(\tau_+) (G_s(\tau_+, \tau, l)G(\tau_+, t, 0) - G(\tau_+, \tau, l)G_s(\tau_+, t, 0))\\
- p(\tau_-) (G_s(\tau_-, \tau, l)G(\tau_-, t, 0) - G(\tau_-, \tau, l)G_s(\tau_-, t, 0)) \} \\
- \{p(\tau_+)(G_s(t_+, \tau, l)G(t_+, t, 0) - G(t_+, \tau, l)G_s(t_+, t, 0)) \\
- p(t_-) (G_s(t_-, \tau, l)G(t_-, t, 0) - G(t_-, \tau, l)G_s(t_-, t, 0)) \}$

\bigskip

ここで $\tau_- \to \tau, \ \tau_+ \to \tau, \ t_- \to t, \ t_+ \to t$ とすれば

\bigskip

$p(\tau_-) \to p(\tau), \ p(\tau_+) \to p(\tau), \ p(\tau_-) \to p(t), \ p(t_-) \to p(t),$

$
G(\tau_-, t, 0) \to G(\tau, t, 0), \ G(\tau_+, t, 0) \to G(\tau, t, 0), \\
G(t_-, t, 0) \to G(t, t, 0), \ G(t_+, t, 0) \to G(t, t, 0), \\
G(\tau_-, \tau, l) \to G(\tau, \tau, l), \ G(\tau_+, \tau, l) \to G(\tau, \tau, l), \\
G(t_-, \tau, l) \to G(t, \tau, l), \ G(t_+, \tau, l) \to G(t, \tau, l), \\
G_s(t_+, \tau, l) - G_s(\tau_-, \tau, l) \to 1/p(\tau), \\
G_s(t_+, t, 0) - G_s(t_-, t, 0) \to 1/p(t).
$

\bigskip

さらに

$
AG(a,t,0) + A'G_s(a,t,0) = 0, \\
AG(a,\tau,l) + A'G_s(a, \tau, l) = 0, \\
BG(b,t,0) + B'G_s(b,t,0) = 0, \\
BG(b,\tau,l) + B'G_s(b,\tau,l) = 0
$
 \ から

$
G_s(a,\tau,l)G(a,t,0) - G(a,\tau,l)G_s(a,t,0) = 0, \\
G_s(b,\tau,l)G(b,t,0) - G(b,\tau,l)G_s(b,t,0) = 0.
$

これらをまとめると, \ 右辺は$-G(\tau,t,0) + G(t, \tau,l)$ に収束する. 

ゆえに $G(\tau,t,0) = G(t, \tau,0), \ G(s,t,0) = G(t,s,0)$ に注意して

$$
G(t, \tau, l)-G(t, \tau, 0)= l \int_{a}^{b} G(t, s, 0) G(s, \tau, l) d s.
$$

これは$l < \tau$ の場合でも同様にして証明できる.



\bigskip
\bigskip

\begin{itembox}[l]{[12]}
$G(t, \tau, \ell)$ を上記問題［11］と同様なグリーン関数とする．$t$ と $\tau$ を固定したとき，$\ell$ を複素数と考えて $(\ell \in \mathbb{C}), G(t, \tau, \ell)$ が $|\ell|<\infty$ において有理型の関数であることを仮定 して，$G(t, \tau, \ell)$ は $\ell=\lambda_{k}$ において 1 位の極をもつことを示せ. ただし，$\lambda_{k}$ は $L$ の$1$つの固有値であるとする.
\end{itembox}

\textbf{解答}

$l = \lambda_k$ の近傍における$G(t,\tau,l)$ のローラン展開を
$$
G(t,\tau,l) = \frac{g_m(t,\tau)}{(l-\lambda_k)^m} + \frac{g_{m-1}(t,\tau)}{(l-\lambda_k)^{m-1}} + \cdots
$$

とすれば, \ 任意の連続関数$f(t)$ に対し

$$
\begin{aligned}
\varphi(t,l) = R_l f(t) &= \int_a^b G(t,\tau,l) f(\tau) d\tau \\
&= \frac{g_m(t)}{(l-\lambda_k)^m} + \frac{g_{m-1}(t)}{(l-\lambda_k)^{m-1}} + \cdots,
\end{aligned}
$$
$$\varphi_n(t) = \displaystyle\int_a^b g_n(t,\tau) f(\tau) d\tau$$.

この両辺に$L-\lambda_k$ を作用させれば, 

$$
(L-\lambda_k)R_l f = (L - l + l - \lambda_k)R_l f = f + (l-\lambda_k)R_l f
$$

に注意して

$$
\left(L-\lambda_{k}\right) \varphi =f+\left(l-\lambda_{k}\right) \varphi
=\frac{\left(L-\lambda_{k}\right) \varphi_{m}}{\left(L-\lambda_{k}\right) m}+\frac{\left(L-\lambda_{k}\right) \varphi^{m-1}}{\left.\left(I-\lambda_{k}\right)\right)_{m-1}}+\cdots
$$

ゆえに

$$
f+\frac{\varphi_{m}}{\left(l-\lambda_{k}\right)^{m-1}}+\frac{\varphi_{m-1}}{\left(l-\lambda_{k}\right)^{m-2}}+\cdots=\frac{\left(L-\lambda_{k}\right) \varphi_{m}}{\left(l-\lambda_{k}\right){ }^{m}}+\frac{\left(L-\lambda_{k}\right) \varphi_{m-1}}{\left(l-\lambda_{k}\right)^{m-1}}+\cdots
$$

この両辺を比べて, \ $m > 1$ ならば

$$
\left(L-\lambda_{k}\right) \varphi_{m}=0, \quad\left(L-\lambda_{k}\right) \varphi_{m-1}=\varphi_{m},
$$

$m=1$ ならば

$$
\left(L-\lambda_{k}\right) \varphi_{1}=0, \quad\left(L-\lambda_{k}\right) \varphi_{0} = f+\varphi_{1}
$$

を得る. よって,

$$
\begin{aligned}
\left\|\varphi_{m}\right\|^{2} & =\left(\varphi_{m}, \varphi_{m}\right)=\left(\varphi_{m}, \left(L-\lambda_{k}\right) \varphi_{m-1}\right) \\
& =\left(\left(L-\lambda_{k}\right) \varphi_{m}, \varphi_{m-1}\right)=0 .
\end{aligned}
$$

したがって $m>1$ に対して$\displaystyle\int_a^b (\varphi_m(t))^2 dt = 0$. $\varphi_m$ は連続関数であるから $\varphi_m = 0$. ゆえに, \ 任意の連続関数$f$ に対して

$$
\int_{a}^{b} g_m(t,\tau) f(\tau)\tau =0
$$

が成り立つ. 以上より$g_m = 0$ であって, \ $G(t,\tau,l)$ は$l = \lambda_k$ において$1$位の極をもつ.


\bigskip
\bigskip

\begin{itembox}[l]{[13]}
固有値問題


\begin{equation*}
\frac{d^{2} x}{d t^{2}}=\ell x, \quad x(-1)=x(1), \quad x^{\prime}(-1)=x^{\prime}(1) \tag{23}
\end{equation*}


を解け.  
\end{itembox}

\textbf{解答}

一般解は $x(t) = A\cos(kt) + B\sin(kt)$.

境界条件$x(-1) = x(1)$ を適用すると,

$A\cos k - B\sin k = A\cos k + B\sin k$ より, \ $B\sin k = 0 ───\textcircled{\scriptsize 1}$

つぎに, \ 境界条件$x'(-1) = x'(1)$ を適用すると, \ 

$x'(t) = -A\sin(kt) + Bk\cos(kt)$ より $Ak\sin k + Bk\cos k = -Ak\sin k + Bk\cos k$

よって, \ $Ak\sin k = 0 ───\textcircled{\scriptsize 2}$

$\textcircled{\scriptsize 1}, \ \textcircled{\scriptsize 2}$ より固有値は$l = (n\pi)^2$, \ $k=n\pi$ に対応する固有関数は$\underline{x(t) = A\cos(n\pi t) + B\sin(n\pi t)}$.



\bigskip
\bigskip

\begin{itembox}[l]{[14]}
偏微分方程式


\begin{equation*}
\frac{\partial^{2}}{\partial t^{2}} u(t, x)=\frac{\partial^{2}}{\partial x^{2}} u(t, x), \quad 0<t<\infty, \quad 0<x<1 \tag{24}
\end{equation*}


をみたし，境界条件 $u(t, 0)=u(t, 1)=0$ をみたし，さらに条件

\[
u(0, x)=\left\{\begin{array}{rl}
x, & 0 \leqslant x \leqslant 1 / 2,  \tag{25}\\
1-x, & 1 / 2 \leqslant x \leqslant 1,
\end{array} \quad \frac{\partial}{\partial t} u(0, x)=0\right.
\]

をみたす解 $u=u(t, x)$ を求めよ.  
\end{itembox}

\textbf{解答}

$u=z(t)w(x)$とおくと

$$
\frac{z^{\prime \prime}}{z}=\frac{w^{\prime \prime}}{w}=l \text {, }
$$
$$
w=c_{1} \cos \sqrt{-l}x + c_{2} \sin \sqrt{-l}x.
$$
$u(t, 0)=u(t, 1)=0$ から $w(0)=w(1)=0$ ．ゆえに
$c_1 = 0, \ \sin\sqrt{-l} = 0$.

以上より$l=-n^2\pi^2$ が得られて

$$
w_{n}=\sin n \pi x,
$$
$$
z^{\prime \prime}=-n^{2} \pi^{2} z
$$

から$z_n = A_n \cos n\pi t + B_n \sin n\pi t$. そこで
$u = \displaystyle \sum_1^{\infty} (A_n \cos n\pi t + B_n \sin n\pi t) \sin n\pi x$ とおくと

$u(0, x)= \displaystyle \sum_{1}^{\infty} A_{n} \sin n \pi x, \quad u_{1}(0, x)=\sum_{1}^{\infty} n \pi B_{n} \sin n \pi x$.\\
$u_{1}(0,x)=0$ から $B_{n}=0$ ．

$u(0,x) = \begin{cases}
x & 0 \le x \le 1/2 \\
1-x & 1/2 \le x \le 1,
\end{cases}$
のフーリエ展開を求めると

$$
A_{n}=\frac{4}{n^{2} \pi^{2}} \sin \frac{n \pi}{2} .
$$

ゆえに解は

$$
u= \underline{\frac{4}{\pi^{2}} \sum_{1}^{\infty} \frac{1}{n^{2}} \sin \frac{n \pi}{2} \cos n \pi t \sin n \pi x} .
$$

\bigskip
\bigskip

\begin{itembox}[l]{[15]}
$J$ を $\mathbb{R}$ のある区間とし，実数値関数 $p(x), q(x)$ はともに $p, q \in C(J)$ であるとす る．複素数値関数 $y=u(x)+i v(x)$ が微分方程式


\begin{equation*}
L y=y^{\prime \prime}+p(x) y^{\prime}+q(x) y=0 \tag{26}
\end{equation*}


の解であるならば，実数値関数 $u(x)$ および $v(x)$ はいずれも方程式 $L y=0$ の解であることを示せ.  
\end{itembox}

\textbf{解答}

$y = u + iv$ を元の方程式に代入すると,

$$
(u +iv)'' + p(x)(u+iv)' + q(x)(u+iv) = 0
$$

よって, \ $\uwave{(u'' + p(x)u' + q(x)u)} + i\uwave{(v'' + p(x)v' + q(x)v)} = 0$

波線部はともに実数値関数であるから, \ $u'' + p(x)u' + q(x)u = v'' + p(x)v' + q(x)v = 0$.

以上より,\ $Lu = Lv = 0$ となるため,

$u(x), v(x)$ が $Ly=0$ の解であることが示された.


\bigskip
\bigskip

\begin{itembox}[l]{[16]}
境界値問題


\begin{equation*}
y^{\prime \prime}=-f(x), \quad y(0)=0, \quad y(1)=0, \quad(f \in C[0,1]) \tag{27}
\end{equation*}


を考える. 対応するグリーン関数を構成し，それを用いて上記境界値問題の解 $y(x)$ を求 めよ.  
\end{itembox}

\textbf{解答}

境界値問題：$y^{\prime \prime} = -f(x), y(0)=0, y(1)=0, f \in C[0,1]$ に対応するグリーン関数をまず構成して一般解$y(x)$ を求める.\\
微分方程式 $\frac{d^2}{d x^2} G(x, \xi)=-\delta(x-\xi)$ より,

$$
G(x, \xi)= \begin{cases}A x & (x<\xi) \\ B(1-x) & (x>\xi)\end{cases}
$$

境界条件 $G(0, \xi)=0$ を適用すると，$A(0)=0$,
$G(1, \xi)=0 \text { を適用すると, } B(1-1)=0$
これらは明らか.


次に，$G(x, \xi)$ の $x=\xi$ での連続性から，$A \xi=B(1-\xi)$ i．e．$A=\frac{B(1-\xi)}{\xi} ───(*)$\\
さらに，1階微分 $\frac{\partial G}{\partial x}$ の $x=\xi$ での不連続性から，


\begin{align*}
& \left.\frac{\partial G}{\partial x}\right|_{x=\xi^{+}}-\left.\frac{\partial G}{\partial x}\right|_{x=\xi^{-}}=-1  \tag{*}\\
& \therefore\left\{\begin{array}{ll}
\frac{\partial G}{\partial x}=A & (x<\xi) \\
\frac{\partial G}{\partial x}=-B & (x>\xi)
\end{array} \quad \text { より, } A+B=1\right.
\end{align*}


これを(*)に代入して解くと，$\frac{B(1-\xi)}{\xi}+B=1 \quad$ より,

$$
B=\xi, A=1-\xi
$$

$\therefore$ グリーン関数は $G(x, \xi)= \begin{cases}(1-\xi) x & (0 \leqq x \leqq \xi) \\ \xi(1-x) & (\xi \leqq x \leqq 1)\end{cases}$

以上より，一般解は

$$
\begin{aligned}
y(x) & =\int_{0}^{1} G(x, \xi) f(\xi) d \xi \\
& =\int_{0}^{x}(1-\xi) x f(\xi) d \xi+\int_{x}^{1} \xi(1-x) f(\xi) d \xi \\
& =\underline{x \int_{0}^{x}(1-\xi) f(\xi) d \xi+(1-x) \int_{x}^{1} \xi f(\xi) d \xi}
\end{aligned}
$$


\bigskip
\bigskip

\begin{itembox}[l]{[17]}
境界値問題


\begin{equation*}
y^{\prime \prime}=-f(x), \quad y(0)=1, \quad y(1)=2, \quad(f \in C[0,1]) \tag{28}
\end{equation*}


を考える。\\
（i），対応するグリーン関数を構成し，それを用いて上記境界値問題の解 $y(x)$ を求めよ．\\
（ii）とくに $f(x) \equiv 1$ のとき，解 $y(x)$ を具体的に書き表せ.  
\end{itembox}

\textbf{解答}

境界問題：$y^{\prime \prime}=-f(x), y(0)=1, \quad y(1)=2, f \in C[0,1]$\\
（i）グリーン関数 $G(x, \xi)$ は，微分方程式 $\frac{\partial^{2}}{\partial x^{2}} G(x, \xi)=-\delta(x-\xi)$\\
を満たすことから，$G(x, \xi)= \begin{cases}A x+B & (x<\xi) \\ C x+D & (x>\xi)\end{cases}$\\
境界条件 $G(0, \xi)=0$ より，$B=0 \quad \therefore G(x, \xi)=Ax$

 \ \ \ \ $G(1, \xi) = 0$ より, $D = -C \quad \therefore G(x, \xi) = C(x-1)$


\begin{equation*}
G(1, \xi)=0 \quad \quad p=-c \quad \therefore G(x, \xi)=c(x-1) \tag{1}
\end{equation*}


$G(x, \xi) の x=\xi$ での連続性より，$A \xi=C(\xi-1) ───\textcircled{\scriptsize 1}$\\
さらに，微分の不連続性より，


\begin{equation*}
\left.\frac{\partial G}{\partial x}\right|_{x=\xi^{+}}-\left.\frac{\partial G}{\partial x}\right|_{x=\xi^{-}}=-1 \Leftrightarrow C-A=-1 ───\textcircled{\scriptsize 2} \tag{2}
\end{equation*}


\textcircled{\scriptsize 1}, \ \textcircled{\scriptsize 2}より，$A=\frac{1-\xi}{2 \xi-1}, C=\frac{\xi}{2 \xi-1}$

$$
\therefore \text { グリーン関数は } G(x, \xi)=\left\{\begin{array}{ll}
\frac{1-\xi}{1-2 \xi} x & (0 \leqq x \leqq \xi) \\
\frac{\xi}{1-2 \xi}(x-1) & (\xi \leqq x \leqq 1)
\end{array} \quad\right. \text { であり, }
$$

一般解は $y(x)=(1-x)+2 x+\int_{0}^{1} G(x, \xi) f(\xi) d \xi$

$$
=1+x+\int_{0}^{x} \frac{1-\xi}{1-2 \xi} x f(\xi) d \xi+\int_{x}^{1} \frac{\xi}{1-2 \xi}(x-1) f(\xi) d \xi
$$

（ii）$f(x) \equiv 1$ のとき,\\
$y^{\prime \prime}=-1, \ y(0)=1, \ y(1)=2$ であるから，\\
$y^{\prime}(x)=-x+C_{1}$\\
$y(x)=-\frac{x^{2}}{2}+c_{1} x+c_{2} \quad\left(c_{1}, c_{2}\right.$ ：定数 $)$\\
境界条件を適用すると，$C_{1}=\frac{1}{2}, \ C_{2}=1$ となるから，

$$
y(x)=-\frac{x^{2}}{2}+\frac{x}{2}+1
$$


\bigskip
\bigskip

\begin{itembox}[l]{[18]}
$I=[a, b]$ とする．$p(x) \in C^{1}(I), p(x)>0$ ，かつ $q(x) \in C(I)$ とする。微分作用素 $L$ を


\begin{equation*}
L y=\left\{p(x) y^{\prime}\right\}^{\prime}+q(x) y=0, \quad(a<x<b) \tag{29}
\end{equation*}


とし，境界条件を

\[
\left\{\begin{array}{l}
B_{a}[y]=A y(a)+A^{\prime} y^{\prime}(a)=0  \tag{30}\\
B_{b}[y]=B y(b)+B^{\prime} y^{\prime}(b)=0
\end{array}\right.
\]

とする。ただし，$A^{2}+A^{\prime 2} \neq 0$ かつ $B^{2}+{B^{\prime}}^{2} \neq 0$ である。任意の $\xi(a \leqslant \xi \leqslant b)$ に対し て，微分方程式


\begin{equation*}
L y=y_{0}(x) y_{0}(\xi) \tag{31}
\end{equation*}


の解 $y(x) \not \equiv 0$ で，境界条件（30）を満足するものは存在しないことを示せ．  
\end{itembox}

\textbf{解答}

微分作用素$L$を持つ方程式 $L y=y_{0}(x) y_{0}(\xi)$ は，左辺が$x$のみの関数，右辺が$x$ と $\xi$ に分離された形.\\
分雖定数法より，$\frac{L\left[y_{0}(x)\right]}{y_{0}(x)}=\frac{L\left[y_{0}(\xi)\right]}{y_{0}(\xi)}=\lambda \quad(\lambda:$ 分離定数 $)$\\
（i）$y_{0}(x) $ 関する方程式\\
$\left\{p(x) y_{0}^{\prime}(x)\right\}^{\prime}+q(x) y_{0}(x)=\lambda y_{0}(x)$ は固有値問題\\
境界条件：$A y_{0}(a)+A^{\prime} y_{0}^{\prime}(a)=0, \ B y(b)+B^{\prime} y_{0}^{\prime}(b)=0$\\
（ii）$y_0(\xi)$ に関する方程式\\
同様に, \ $y_0(\xi)$ に対しても$\{p(\xi)y_0'(\xi) \}' + q(\xi)y_0(\xi) = \lambda y_0(\xi)$

分離定数法より, \ (i), (ii) の$2$ つの方程式は互いに独立して成立.

$2$ つの変数$x$ と $\xi$ が独立しているため, \ 境界条件のもとで$2$ つの独立した固有値問題が一致しない. $\therefore Ly=y_{0}(x) y_{0}(\xi)$ を満たす解は存在しない.

\bigskip
\bigskip

\begin{itembox}[l]{[19]}
境界値問題


\begin{equation*}
y^{\prime \prime}+y=-f(x), \quad y(0)=y(\pi)=0, \quad(f \in C[0, \pi]) \tag{32}
\end{equation*}


を考える。斉次境界値問題の解 $y_{0}(x)=\sqrt{2 / \pi} \sin x, y_{0}(x)$ と 1 次独立な斉次境界値問題 の解 $y_{1}(x)$ ，非斉次境界値問題の特殊解


\begin{equation*}
w(x, \xi)=\frac{1}{2 \pi} \sin \xi \sin x-\frac{1}{\pi} x \cos x \sin \xi \tag{33}
\end{equation*}


に対して，範囲 $0 \leqslant x \leqslant \xi$ に対応するグリーン関数 $G(x, \xi)$ を具体的に求めよ．  
\end{itembox}

\textbf{解答}

グリーン関数$G(x, \xi)$ は以下を満たす.

$G''(x, \xi) + G(x, \xi) = \delta(x - \xi)$ \ \ 境界条件：$G(0,\xi) = G(\pi, \xi) = 0$

$\underline{G(x, \xi) = \begin{cases}
\frac{1}{2\pi} \sin\xi\sin x - \frac{1}{\pi}x\cos x\sin\xi & (0\le x\le\xi) \\
\frac{1}{\pi}(\xi - \pi) \cos\xi\sin x & (\xi\le x \le\pi)
\end{cases}}$

\bigskip
\bigskip

\begin{itembox}[l]{[20]}
条件\\
（a）$G(x, \xi)$ は $a \leqslant x, \xi \leqslant b$ で連続で，かつ任意の $\xi$ に対して $x \neq \xi$ で 2 回連続微分可能である。\\
（b）関係式


\begin{equation*}
\left.\lim _{\varepsilon \rightarrow 0} G_{x}^{\prime}(x, \xi)\right|_{\xi=x-\varepsilon} ^{\xi=x+\varepsilon}=G_{x}^{\prime}(x, x+0)-G_{x}^{\prime}(x, x-0)=\frac{1}{p(x)} \tag{34}
\end{equation*}


が成り立つ。\\
（c）任意の $\xi$ に対して，$G(x, \xi)$ は $\quad x \neq \xi$ で微分方程式


\begin{equation*}
\frac{d}{d x}\left\{p(x) \frac{d G}{d x}(x, \xi)\right\}+q(x) G(x, \xi)=0 \tag{35}
\end{equation*}


をみたす。\\
（d）任意の $\xi$ に対して，$G(x, \xi)$ は境界条件


\begin{equation*}
A G(a, \xi)+A^{\prime} G_{x}^{\prime}(a, \xi)=0, \quad B G(b, \xi)+B^{\prime} G_{x}^{\prime}(b, \xi)=0 \tag{36}
\end{equation*}


をみたす。\\
上述のような条件（a），（b），（c），（d）を満足する関数 $G(x, \xi)$ は対称性をもつことを示せ．  
\end{itembox}

\textbf{解答}

以下を証明する：
\[
G(x, \xi) = G(\xi, x)
\]


2つの関数 \( G(x, \xi) \) および \( G(\xi, x) \) に対して, \ グリーンの第二定理を適用すると,

\[
\int_a^b \left[ G(x, \xi) \frac{d}{d\xi} \left\{ p(\xi) \frac{dG}{d\xi}(x, \xi) \right\} - G(\xi, x) \frac{d}{dx} \left\{ p(x) \frac{dG}{dx}(x, \xi) \right\} \right] dx = 0
\]

積分範囲内の微分方程式を用いることで, \ 上記の積分は境界項に帰着する：

\[
\left[ p(x) \left( G(x, \xi) \frac{dG}{dx}(x, \xi) - G(\xi, x) \frac{dG}{d\xi}(x, \xi) \right) \right]_a^b = 0
\]

境界条件 (d) により, \ これらの境界項はゼロになる.

次に, \ 条件(b) の微分の不連続性より
\[
G'_x(x, \xi+0) - G'_x(x, \xi-0) = \frac{1}{p(x)}
\]

対称的に書き換えると
\[
G'_x(\xi, x+0) - G'_x(\xi, x-0) = \frac{1}{p(\xi)}
\]

これにより、対称性
\[
G(x, \xi) = G(\xi, x)
\]

が成立する.



\bigskip
\bigskip

\begin{itembox}[l]{[21]}
グリーン関数 $G(x, \xi)$ を求めて，次の境界値問題を解け．


\begin{equation*}
y^{\prime \prime}+y=-f(x), \quad y(0)=y\left(\frac{\pi}{2}\right)=0, \quad\left(f \in C\left[0, \frac{\pi}{2}\right]\right) \tag{37}
\end{equation*}  
\end{itembox}

\textbf{解答}

与えられた境界値問題を解くために, \ 次の方程式を満たすようなグリーン関数 \( G(x, \xi) \) を求める.


\[
\frac{d^2}{dx^2} G(x, \xi) + G(x, \xi) = \delta(x - \xi),
\]

および境界条件：
\[
G(0, \xi) = 0, \quad G\left(\frac{\pi}{2}, \xi\right) = 0.
\]

(i) \( x < \xi \) の場合

一般解は次の形になる：
\[
G_1(x, \xi) = A \sin(x) + B \cos(x).
\]

境界条件 \( G(0, \xi) = 0 \) より
\[
B = 0.
\]

したがって：
\[
G_1(x, \xi) = A \sin(x).
\]

(ii) \( x > \xi \) の場合

同様に、一般解は次の形になる：
\[
G_2(x, \xi) = C \sin(x) + D \cos(x).
\]

境界条件 \( G\left(\frac{\pi}{2}, \xi\right) = 0 \) より
\[
C = 0.
\]

したがって
\[
G_2(x, \xi) = D \cos(x).
\]


デルタ関数の性質より, \ 不連続条件が存在する

\[
G'_x(x=\xi+0) - G'_x(x=\xi-0) = 1.
\]

それぞれの導関数を計算すると
\[
\frac{d}{dx} G_1(x, \xi) = A \cos(x),
\]
\[
\frac{d}{dx} G_2(x, \xi) = -D \sin(x).
\]

したがって, \ 不連続性の条件は次のようになる
\[
-A \cos(\xi) - (-D \sin(\xi)) = 1.
\]

\[
-A \cos(\xi) + D \sin(\xi) = 1.
\]

また, \ 連続性条件
\[
A \sin(\xi) = D \cos(\xi).
\]

これを連立して解くと
\[
A = \frac{\cos(\xi)}{\sin(\xi)}, \quad D = \frac{\sin(\xi)}{\cos(\xi)}.
\]


以上の結果をまとめると、グリーン関数は次のように表される：

\[
\begin{aligned}
G(x, \xi) &=
\begin{cases}
\frac{\sin(x) \sin\left(\frac{\pi}{2} - \xi\right)}{\sin\left(\frac{\pi}{2}\right)}, & x < \xi, \\
\frac{\sin(\xi) \sin\left(\frac{\pi}{2} - x\right)}{\sin\left(\frac{\pi}{2}\right)}, & x > \xi.
\end{cases}  \\
&= \underline{\begin{cases}
\sin(x) \sin\left(\frac{\pi}{2} - \xi\right), & x < \xi, \\
\sin(\xi) \sin\left(\frac{\pi}{2} - x\right), & x > \xi.
\end{cases}}
\end{aligned}
\]


これにより、与えられた境界値問題が解かれた.


\end{document}
